% =====================================================================
\section{Evaluation method}
\label{sec:method}
% =====================================================================

Every comparison in Sections~\ref{sec:chain}--\ref{sec:width} ran in a
discrete-time simulator of the gateway (the \emph{plant}), with one
controller configuration (an \emph{arm}) per run. A generated
specification in the artifact lists every rate, shape, mix, budget,
target and constant; this section gives what is needed to read the
results.

\subsection{Plant}
\label{sec:plant}

\begin{table}[t]
\caption{Model tiers of the plant: price per cache miss, base latency per
AI request kind, and answer quality (mean of 11 benchmark scores of the
Qwen2.5 instruction-tuned models~\cite{qwen25}, divided by 100; priced only
in the QUALITY campaign). The prices form a synthetic 1:10:100 ladder;
our own GPU bench measured 1:1.475:1.518 for the three models.}
\label{tab:tiers}
\centering\tablefont
% generated by make_assets.py -- do not edit
\begin{tabular}{@{}l c r r r c@{}}
\toprule
 & Price & \multicolumn{3}{c}{Base latency (ms)} & \\
\cmidrule(lr){3-5}
Tier (model) & (\$/miss) & Chat & Embed & Agent & Quality \\
\midrule
none & 0 & \multicolumn{3}{c}{30,000, no model call} & 0.000 \\
small (0.5B) & $10^{-4}$ & 300 & 80 & 6,000 & 0.284 \\
mid (3B) & $10^{-3}$ & 800 & 150 & 2,500 & 0.548 \\
large (7B) & $10^{-2}$ & 2,000 & 400 & 3,500 & 0.641 \\
\bottomrule
\end{tabular}

\end{table}

A replica serves 100 work units per second; a CRUD read costs 1 unit, a
write 2, an embedding 5, a chat request 20 and an agent step 40, so AI
traffic loads the replicas too. The semantic cache is a synthetic hit
curve: $c$\,MB answer a fraction $h(c)=0.85\,c/(c+256)$ of a request kind's
cacheable share (40\% of chat, 70\% of embedding and 10\% of agent
requests), with no model of prompt similarity, eviction or drift. Fitted to
200{,}000 real LMSYS-Chat-1M turns, the same form gives the much lower
$0.285\,c/(c+19)$, which Section~\ref{sec:coordination} tests. Models run
outside the plant; Table~\ref{tab:tiers} gives each tier's price per cache
miss, base latency and quality. An agent request on the small tier takes
6\,s at base, longer than on the mid (2.5\,s) or the large tier (3.5\,s);
this non-monotone ladder is a modelling assumption (tool-loop retries on a
weak model) that we did not measure, and it makes the cheapest tier miss
the agent targets of premium and standard tenants.
Congestion multiplies latency by $g(\rho)=1/(1-\rho)$, where $\rho$ is the
work left after cache hits over the replicas' capacity; past $\rho=0.95$ the
factor grows quadratically and is capped at 50. CRUD p95 latency is
$1.4\,\ell_0\,g(\rho)$ for a class base latency $\ell_0$; with $q_k$ the
traffic share and $h_k$ the hit rate of AI kind $k$, AI p95 latency (ms) is
\begin{equation}
\label{eq:lat}
\ell^{\mathrm{AI}}=1.4\sum_k q_k\bigl[20\,h_k+(1-h_k)\,L_{\tau,k}\,g(\rho)\bigr],
\end{equation}
with $L_{\tau,k}$ the tier's base latency. SLO targets of 150\,ms (CRUD) and
2.5\,s (AI) are scaled by the tenant's class: 1 for premium, 2.5 for
standard and 8 for best-effort. On a shed step, CRUD and cache hits are
served as usual, but each AI cache miss gets no model call and enters
(\ref{eq:lat}) at a flat 30\,s; we count these misses as \emph{shed
requests}. Replicas cost \$0.048 per replica-hour and cache memory \$0.005
per GB-hour. Actuation is instantaneous except in the ACTUATION campaign,
where replica changes take effect one interval late and enlarged caches
start half-warm; tier switches never lag.

\emph{What dominates the bill.} One replica-step costs
$\$1.3\times10^{-4}$, about the price of one small-tier cache miss; a
large-tier miss costs $\$10^{-2}$. A standard tenant's hourly budget of \$5
allows \$0.0139 per step, while a cacheable-AI tenant at its base rate sends
about 50 cache misses per step: \$0.005 on the small tier, \$0.05 on the
mid tier. Across the main matrix, model spend was 77\% of \jcac{}'s bill,
replica time 23\% and cache memory 0.2\%. A per-tenant budget can therefore
be kept only through the cache and tier knobs or by shedding; the replica
knob moves mostly latency. Agentic tenants meet their premium and standard
targets only on the mid tier or above, which no budget class pays for during
their bursts (about 50 mid-tier misses, or \$0.05, per step); best-effort
tenants meet theirs on the small tier.

\emph{What is synthetic.} The price ladder of Table~\ref{tab:tiers}
(1:10:100) is an assumption modelled on market per-request prices; every
campaign here uses it except QUALITY, which uses the measured 1:1.475:1.518.
The non-monotone agent latencies are an assumption too, and our own
measurements contradict them: a GPU bench measured the three models at
2.1, 3.1 and 3.2\,s per request, and the live-fitted plant is fastest on the
small tier for every request kind. Congestion multiplies the external
model's latency, although a gateway pod's queue would add to it, so
replicas affect AI latency more than a pod's service time warrants. A cache
hit costs 20\,ms and nothing else, and AI p95 latency is 1.4 times a
traffic-weighted mean, so hits are credited linearly.

\subsection{Workloads and design matrix}
\label{sec:workloads}

The main matrix crosses five workload classes (steady and bursty CRUD,
cacheable and uncacheable AI, and agentic) with four eight-tenant mixes
(uniform; premium-heavy; best-effort-heavy; and a whale mix in which one
tenant sends four times the base demand) and three cluster sizes. The sizes cap each
tenant at 6, 10 or 16 replicas and the tenants together at 24, 48 or 96
replicas and 2, 4 or 8\,GB of cache, below the sum of the individual caps,
so tenants contend for one pool. Hourly budgets are \$2, \$5 and \$8 for
best-effort, standard and premium tenants (\$12 for the whale). This gives 60 \emph{design points}, each run for 120
ten-second steps (20 simulated minutes); every arm runs five
\emph{replications} per design point, one per seed. In the campaigns
reported here all arms share seeds and therefore see identical demand; the
pre-audit campaigns of Section~\ref{sec:chain} did not, because their seeds
included the arm's name. Demand shapes are periodic (sine, sawtooth and
square-wave bursts with per-tenant phases).

Four other substrates test generalisation. HELDOUT runs two
Markov-modulated workloads that appear in no tuning step or earlier
campaign: demand switches between 0.4 and 2.5 times its base rate, upward
with probability 0.08 and downward with 0.25 per step, under lognormal
noise ($\sigma=0.15$), on 24 design points. TRACE\_AUDITED (TA) replays 96
six-hour windows of the real BurstGPT trace~\cite{burstgpt} as chat demand
for the uniform mix of eight standard tenants on the medium cluster, one
replication per window. A \emph{live-fitted} plant
follows measurements of the deployed gateway: a replica serves at least
1{,}123 units per second (no measured run reached congestion) and AI
requests never load it. WIDTH adds a \emph{coupled} plant in which AI work
loads the replicas three times as heavily and the mid and large tiers cost
two and four times the small tier's replica work.

\subsection{Comparators}
\label{sec:comparators}

An HPA arm sets $\lceil r\rho/\rho^\ast\rceil$ replicas for utilisation
$\rho$ and target $\rho^\ast$, following the Kubernetes rule~\cite{k8shpa};
a KEDA arm sets $\lceil q/q^\ast\rceil$ for request rate $q$ and
per-replica target $q^\ast$~\cite{keda}. Both clamp replica changes to
$\pm2$ per interval, as \jcac{} does, and keep a fixed cache and tier. Their
targets were tuned on a separate slice (three of the five workloads, two
of the four mixes, two replications) over grids whose optima are interior,
then re-tuned per cluster size; the \emph{capped} variants also obey the
cluster caps that bind \jcac{}. The tuned utilisation targets settled at
5--10\% (the Kubernetes default is 80\%), because replicas are cheap in
this plant and $J$ is nearly flat across low targets (about 1\% from 0.01
to 0.1); one target per cluster size, pooled over workloads, therefore
over-provisions the CRUD classes. A \emph{layered stack} runs an HPA rule, a
cache policy and a tier rule independently. We test four tier rules: the
registered rule (rank tiers by latency, ignore price); a cost-aware rule
(the cheapest tier expected to meet the latency target); a pin to the
cheapest model for all traffic (the \emph{cheapest-tier stack}); and, in QUALITY, a quality-aware router. A static equal share of
the cluster on the cheapest tier is the null controller.

The capped HPA and KEDA arms and the quality-aware stack obey the cluster
caps; the other layered stacks were registered without them. Auditing the
runs for this paper, we found that the cheapest-tier stack used up to twice
the small cluster's replica cap (about 48 replicas against 24, on every
step) and exceeded the cache cap on most steps of the small and medium
clusters, while \jcac{} never exceeded either cap. We report its registered
comparisons as scored and, beside them, an exploratory re-run of the same
stack held to the caps by the wrapper that caps HPA and KEDA, with its
registered tuning (it was not re-tuned under the caps).

Two asymmetries remain by design and are measured rather than removed.
Only \jcac{} can shed, and only \jcac{} enforces the per-tenant budget. The
BUDGET\_MATCHED campaign (BM) re-ran capped HPA, capped KEDA and the
cheapest-tier stack under \jcac{}'s budget filter on the registered seeds.
The filter can only trim their replicas, because their cache and tier rules
stay their own; since replicas are a small part of the bill
(Section~\ref{sec:plant}), these arms remained over budget on 2.2--3.2\%
of tenant-steps (\jcac{}: 0.9\%), and BM measures what the filter does to
their latency, not a budget-compliant comparator. No arm in this study can
both serve every request and keep every budget.

\subsection{Objective}
\label{sec:objective}

Arms are scored on
\begin{equation}
\label{eq:J}
J=\bar C/c_0+2\,\bar v+0.5\,\bigl(1-\overline{\mathrm{Jain}}\bigr),\qquad
\bar C=C_{\mathrm{run}}/(8\times119),
\end{equation}
with $C_{\mathrm{run}}$ a run's total cost and $\bar v$ and
$\overline{\mathrm{Jain}}$ its means over tenant-steps and steps. $J$ uses
the planner's own weights. It prices violation, whereas (\ref{eq:plan})
plans on the logarithm of excess; a planner variant that plans on
violation instead changed $J$ by $+0.0047$ $[-0.0025,+0.0121]$. The
registered scorers of the seven core campaigns and of RV divide cost by 119
steps, one fewer than the 120 the simulator scores, so $\bar C$ is 0.84\%
above the cost per scored tenant-step; re-scoring with 120 changes one of
79 verdicts (RV-W4, from FAIL to PASS), on which no conclusion rests. TA, CS
and BM divide by each run's own scored steps.

Three properties of $J$ bear on every result below. First, it ignores
budget overruns, which only the comparators incur (3.1--10.4\% of their
tenant-steps). Second, it prices a shed only through violation, which is
traffic-weighted and capped at 1 per request family: shedding all AI
traffic of one tenant for one step raises that step's $v$ by at most the
AI share $q_{\mathrm{AI}}$ of its requests, so it adds at most
$2q_{\mathrm{AI}}$ to the per-step term whose cost part is
$\mathrm{cost}/c_0$. An outage is therefore worth it to $J$ whenever the
step's model spend exceeds $2q_{\mathrm{AI}}c_0=\$0.02\,q_{\mathrm{AI}}$,
about twenty mid-tier misses; for an agentic tenant at the mid tier during
a burst it always is. The bound is often loose: because AI p95 is a scaled
mean, the cache hits a shed tenant still answers pull its latency down, and
where the small tier already misses the target (premium agent traffic) a
shed adds no violation at all. Third, the registered weights are a choice,
and the switching penalty $\lambda=0.05$ per changed knob is a hand-set
constant. With it, \jcac{}'s replicas ratchet upward (11.3 increases and 0.9
decreases per run); with the smaller penalty the code documents as the
corrected value (half a replica-step, 0.0067), replicas move both ways but
\jcac{}'s $J$ worsens by 0.031, more than the equivalence margin
(exploratory). Section~\ref{sec:threats} reports how verdicts move with the
weights.

\subsection{Statistics}
\label{sec:stats}

The unit of analysis is the design point (replications averaged) or the
trace window, and every comparison pairs two arms on the same units and
seeds. A directional claim passes only if a one-sided Wilcoxon signed-rank
test clears its Holm threshold~\cite{holm1979} within its campaign's
registered family \emph{and} the 95\% bootstrap interval of the mean paired
difference (10{,}000 resamples) lies on the registered side of zero. Two
arms \emph{tie} when that interval spans zero; only an equivalence
reading~\cite{lakens2017} bounds a difference: equivalent (EQ) when the
interval lies inside the registered margin of $\pm0.025$ in $J$, the
smallest difference the design detects at $n=60$ with power 0.8 (about 6\%
of \jcac{}'s $J$; HELDOUT, with $n=24$, uses the same margin). A difference
can therefore pass its test and still be practically equivalent. Most
campaigns also carry a non-inferiority guard on the uncapped SLO excess
(one-sided 95\% bound below $+0.05$). Replications within a design point
are strongly correlated (intraclass correlation of paired differences
0.97--0.99 where we computed it), so five replications carry little more
information than one, and the intervals resample design points: they
generalise only over the sampled workloads, mixes and sizes.

\subsection{The level-playing-field campaign}
\label{sec:levelmethod}

A final audit of the eleven campaigns (Sections~\ref{sec:shed}
and~\ref{sec:threats}) found leaks on both sides: a free outage and a budget
filter that only \jcac{} could use, which flattered it; a comparator above
the cluster caps and an over-punitive never-shed fallback, which flattered
the comparators; an unmeasured agent latency row; and a switching penalty
that ratchets \jcac{}'s replicas upward. The twelfth campaign, LEVEL, removes
them together, with every change opt-in so that published arms replay
bit-identically.

\emph{Objective.} $J_L$ is $J$ plus an outage price: every refused AI
request is charged the small tier's price,
\begin{equation}
\label{eq:JL}
J_L = \frac{C_{\mathrm{run}} + n_{\mathrm{shed}}\,p_{\mathrm{small}}}{8\,T\,c_0}
      + 2\bar v + 0.5\bigl(1-\overline{\mathrm{Jain}}\bigr),
\end{equation}
with $T$ the run's scored steps and $n_{\mathrm{shed}}$ recovered exactly
from the stored request counters, so refusing service never saves money
relative to the cheapest tier.

\emph{Comparators.} Every arm obeys the cluster caps. A new
\emph{budget-gated} stack is the cheapest-tier stack whose gateway refuses AI
model calls in any interval the tenant's budget cannot pay for (tier
``none'', cache hits still answered), the way a gateway budget does; it can
shed exactly when \jcac{}'s budget filter would force it to. The gated and
ungated cheapest-tier stacks, capped HPA and capped KEDA were re-tuned per
cluster size under $J_L$ on the tuning slice (no optimum on a grid edge).

\emph{Treatment.} Four variants of \jcac{}, each a documented fix of an
audited defect, were declared before scoring and evaluated under $J_L$ on the
same tuning slice; the declared rule (lowest mean over the three sizes)
chose the variant that pays the outage price in its own objective (the
budget filter still tests real spend), serves on the cheapest affordable
tier before it sheds, uses a switching penalty of half a replica-step, and
forecasts with damped Holt. On that slice the chosen variant still scored
worse than every comparator on the small and medium clusters, and the
protocol predicted failures accordingly.

\emph{Substrates.} Five experiments on fresh seeds (8,832 runs): the main
matrix (primary); the main matrix under a monotone agent latency row (agent
latency 3.125 times chat latency per tier: 0.94, 2.5 and 6.25\,s); the main
matrix with replicas one interval late; the two held-out workloads; and the
96 BurstGPT windows. One Holm family holds eleven tests; coordination is read
against a tighter margin of $\pm0.01$, with the earlier $\pm0.025$ beside it.
The protocol and scorer were pushed to the public repository 15\,s before
the runner started, with the push recorded by a CI run; the predictions
were written after seeing the treatment and the comparators on the tuning
slice, and they are mostly failures.

\subsection{Registration}
\label{sec:prereg}

Each campaign's protocol fixes its arms, design, tests, predictions and
scorer before its first run. Table~\ref{tab:campaigns} lists the twelve
campaigns reported here. The protocols of the seven core campaigns and of
LEVEL were pushed to the public repository before their first runs, 0.2\,min
to 14.9\,h ahead. The protocols of the four follow-ups (RV, TA, CS and BM) and
of the request-counting re-run SA were committed to the authors' local
repository before their runs (0.1\,min to 25.3\,h ahead) and published only
afterwards, so their order rests on commit times we control. Six
protocol--campaign pairs started within two minutes of their protocol's
commit. The campaigns are confirmatory one by one but adaptive as a
sequence: each was designed after an exploratory re-analysis of its
predecessor, mostly on the same 60 design points, and with replications so
strongly correlated a fresh seed adds little independent evidence. Only
HELDOUT ran on design points that no earlier analysis had seen. Two earlier
campaigns cited in Section~\ref{sec:chain} (EVICTION\_PARITY and
ORDER\_PERMUTATION) were committed after their first runs and are disclosed
as such. The 69 $p$-valued tests of the twelve campaigns are listed with
their verdicts in~\ref{app:tests}; the programme-wide ledger of all 63
protocols counts 157 such tests, of which 82 passed and 64 survive a
programme-wide Holm correction, plus 88 interval- or rule-based readings
without a $p$-value.

\begin{table}[t]
\caption{The twelve registered campaigns: seven core campaigns, four
follow-ups and the final level-field campaign (LEVEL, five experiments). Points: design points or BurstGPT windows. PASS: $p$-valued
tests passed, out of those registered (equivalence readings excluded).
BM runs three new arms against registered rows; the SA re-run of the JN
and TA runs (request counters only) is descriptive and not counted.}
\label{tab:campaigns}
\centering\footnotesize\setlength{\tabcolsep}{3.5pt}
% generated by make_tables.py from tab_campaigns_base.tex (paper-ieee make_assets.py) -- do not edit
\begin{tabular}{@{}l l l r r r r@{}}
\toprule
Campaign & Tag & Question & Arms & Points & Runs & PASS \\
\midrule
FAIR\_J & FJ & Fair comparators & 7 & 60 & 2,100 & 4/5 \\
RETUNED & RT & Per-size re-tuning & 6 & 60 & 1,800 & 1/3 \\
JOINTNESS & JN & Coordination, stacks, caps & 12 & 60 & 3,600 & 5/7 \\
ACTUATION & AC & One-interval actuation lag & 6 & 60 & 1,800 & 3/4 \\
QUALITY & QL & Answer quality priced & 8 & 60 & 2,400 & 1/5 \\
HELDOUT & HO & Unseen bursty demand & 7 & 24 & 840 & 0/5 \\
WIDTH & WD & Step width $\times$ coupling & 4 & 3$\times$48 & 2,880 & 5/5 \\
\midrule
Subtotal & & & & & 15,420 & 19/34 \\
\midrule
REVIEWER & RV & Wide steps, live plant & 6+7 & 2$\times$60 & 3,900 & 2/8 \\
TRACE\_AUD. & TA & Real BurstGPT demand & 7 & 96 & 672 & 1/4 \\
CACHE\_SENS. & CS & Three cache-hit curves & 4 & 3$\times$60 & 3,600 & 0/6 \\
BUDGET\_MAT. & BM & Comparators on the budget & 3 & 60+96 & 1,188 & 4/6 \\
\midrule
LEVEL & LV & Level field, final & 8 & 5 subs. & 8,832 & 1/11 \\
\midrule
Total & & & & & 33,612 & 27/69 \\
\bottomrule
\end{tabular}

\end{table}
